\documentclass[platex,dvipdfmx,12pt]{jsarticle}
%使ってるのがplatexだからクラス指定はplatex.間違えてもuplatexにしないように

%==============================================================================
% 要件定義書テンプレート
%   土台 : kosen-report-writer の latex_snipets.json (reportTemplate) のプリアンブル
%   再現 : 要件定義書サンプル.pdf（表紙／改訂履歴／目次／1.概要〜5.補足／以上）
%
% ■ビルド
%   latexmk -outdir=latex-output requirements-template.tex
%   （.latexmkrc により platex → pbibtex → platex → dvipdfmx で処理される）
%
% ■記入ルール（後続チャットで内容を当てはめる際の約束事）
%   ・【...】 が差し替え対象のプレースホルダ．記入したら【】ごと消す
%   ・章(\section)・節(\subsection)・■見出し(\reqlabel)の構成はサンプルに準拠．
%     不要な ■ブロックは丸ごと削除，足りなければ複製してよい
%   ・表の行は「セル & セル \\ \hline」．行は必要なだけ増減する（最終行にも \\ \hline）
%   ・セル内の箇条書きは「・項目1\par・項目2」のように \par で区切る
%   ・本文の箇条書きは reqbullet 環境
%   ・図は assets/ に置き，\reqimage{assets/ファイル名} を差し替える
%     （ファイルが無い間は「※図を添付」の枠が出る．ファイル名に _ は使わない）
%   ・句読点はサンプルに合わせて「、。」
%   ・特殊文字は \% \& \_ \# \$ とエスケープ．丸数字などは \ajMaru{1}（otf）
%   ・ファイル名などをインラインで書くときは \texttt を使わずそのまま書く
%
% ■独自マクロ一覧（プリアンブルで定義）
%   \reqlabel{名前}       ■見出し。直後に \begin{...}（表）が続くときは表の先頭行に組み込まれ，
%                         見出しだけがページ末に取り残されない（段落が続くときは単独の見出し）
%   \reqth{文字}          表ヘッダー用（白・太字）
%   \reqimage[幅]{パス}    画像．無ければ枠表示
%   reqbullet 環境        本文の「・」箇条書き（\item で項目）
%   reqtable 環境         ヘッダー付き表   \begin{reqtable}{列指定}{ヘッダー行}
%   reqlist 環境          ヘッダー無しの2列表（項目名＋値）\begin{reqlist}{列指定}
%   reqcategory 環境      カテゴリー | 詳細 の表（1章）
%   reqitems 環境         項目 | 説明 | 具体的な要件 の表（4章）
%   列指定 : K{比率}=太字ゴシックの項目列  P{比率}=通常列  M{比率}=中央揃え列
%            比率は表幅に対する割合で，合計が 1 になるようにする（例 K{.3}P{.7}）
%==============================================================================

\usepackage{graphicx}
\usepackage[hyphens]{url}
\usepackage{otf}
\usepackage{multirow}
\usepackage{amsmath}
\usepackage{here}
\usepackage{geometry}
\usepackage{ascmac}

\usepackage[hang,small,bf]{caption}
\usepackage[subrefformat=parens]{subcaption}
\captionsetup{compatibility=false}
\usepackage{textcomp}

\usepackage{listings} %日本語のコメントアウトをする場合jvlisting（もしくはjlisting）が必要
%ここからソースコードの表示に関する設定
\lstset{
  basicstyle={\ttfamily},              % 基本:等幅フォント
  identifierstyle={\small},            % 識別子
  commentstyle={\smallitshape},        % コメント:斜体
  keywordstyle={\small\bfseries},     % キーワード:太字
  stringstyle={\small\ttfamily},      % 文字列
  frame={tb},                           % 上部と下部に線を表示
  breaklines=true,                      % 行が長い場合に折り返す
  columns=[l]{fullflexible},            % 列幅を自動調整する（見た目が良くなる）
  numbers=left,                         % 行番号を左側に表示
  xrightmargin=0zw,                     % 右マージンのサイズ．
  xleftmargin=1.6zw,                    % 左マージンのサイズ．行番号が2桁でも行左端からはみ出ない値．
  numberstyle={\scriptsize},           % 行番号のスタイル．スクリプトサイズのフォントを使用
  stepnumber=1,                         % 何行ごとに行番号を表示するか
  numbersep=1zw,                        % 行番号とソースコードの間の距離
  lineskip=-1.4ex,                      % ソースコードの行間（結構詰め気味）
  morestring=[b]",                     %日本語等のStringを正常にダブルクオーテーションで囲む
  showstringspaces=false,               %Stringに囲まれた部分ではスペースがSymbol付きで表示されるのでそれを消す
}
\renewcommand{\lstlistingname}{ソースコード}

\geometry{left=20mm,right=20mm,top=20mm,bottom=30mm}
\renewcommand{\baselinestretch}{1.2}

%------------------------------------------------------------------------------
% ここから要件定義書用の追加（上までがスキル標準のプリアンブル）
%   longtable : ページをまたぐ長い表（ヘッダー行を各ページに再掲）
%   needspace : 見出し・■見出しだけがページ末尾に取り残されるのを防ぐ
%   xcolor[table] : サンプルの水色ヘッダー・縞模様・青枠のための色指定
%------------------------------------------------------------------------------
\usepackage{longtable}
\usepackage{needspace}
\usepackage[table]{xcolor}

% サンプルPDFの配色
\definecolor{reqblue}{RGB}{68,114,196}   % 表紙の枠・見出しの青
\definecolor{reqbox}{RGB}{217,226,243}   % 章見出しの塗り
\definecolor{reqhead}{RGB}{91,155,213}   % 表ヘッダー
\definecolor{reqlight}{RGB}{222,234,246} % 表の縞（薄い水色）
\definecolor{reqrule}{RGB}{157,195,230}  % 表の罫線

\makeatletter %sectionをchapterっぽくする.
% 見出し番号を「1.」「(1)」形式にする（サンプルPDFに合わせた設定）．
% スキル標準の「第1章」形式にする場合は pre=第 / post=章 に戻し，\thesubsection の行を削除する
\renewcommand{\presectionname}{}
\renewcommand{\postsectionname}{.}%はコメントですがつけ忘れると\section*でバグります
\renewcommand{\thesubsection}{(\arabic{subsection})}

% 図番号を1.5のようにする
\renewcommand{\theequation}{\arabic{section}.\arabic{equation}}
\renewcommand{\thefigure}{\arabic{section}.\arabic{figure}}
\renewcommand{\thetable}{\arabic{section}.\arabic{table}}
% jsarticle は章ごとにカウンタがリセットされないため追加（無いと2章の最初の図が「図2.2」になる）
\@addtoreset{equation}{section}
\@addtoreset{figure}{section}
\@addtoreset{table}{section}

% 章見出し : 青枠＋薄青塗りのボックス（スタイル末尾のマクロが見出し全体を受け取る）
\newcommand{\reqsectionbox}[1]{%
  \begingroup\setlength{\fboxrule}{1.6pt}\setlength{\fboxsep}{5pt}%
  \noindent\fcolorbox{reqblue}{reqbox}{\parbox{\dimexpr\linewidth-2\fboxsep-2\fboxrule\relax}{\raggedright #1}}%
  \endgroup}
% 節見出し : 左に青いバー，下に青い罫線
\newcommand{\reqsubsectionbox}[1]{%
  \noindent\textcolor{reqblue}{\rule[-.8ex]{.5zw}{3ex}}\hspace{.5zw}%
  \begin{minipage}[t]{\dimexpr\linewidth-1zw\relax}\raggedright #1\end{minipage}\par\nobreak
  \vspace{-.7\baselineskip}\noindent\textcolor{reqblue}{\rule{\linewidth}{1.2pt}}\par}

\renewcommand{\section}{%
\clearpage
\@startsection{section}{1}{\z@}%
{\Cvs \@plus.5\Cdp \@minus.2\Cdp}% 見出し前に付けられる空白量
{.5\Cvs \@plus.3\Cdp}% 見出し後に付けられる垂直方向の空白量
{\normalfont\Large\headfont\reqsectionbox}}% 見出しのスタイル
\renewcommand{\subsection}{%
\Needspace{6\baselineskip}%
\@startsection{subsection}{2}{\z@}%
{\Cvs \@plus.5\Cdp \@minus.2\Cdp}%
{.5\Cvs \@plus.3\Cdp}%
{\normalfont\large\headfont\reqsubsectionbox}}
\makeatother

\renewcommand{\contentsname}{目次}
\setcounter{tocdepth}{2}

\makeatletter
% 表の列数を数える（列指定 K{..}P{..}M{..} の文字数）
\newcount\req@ncols
\newcommand{\req@countcols}[1]{\req@ncols\z@
  \@tfor\req@t:=#1\do{%
    \def\req@a{K}\ifx\req@t\req@a\advance\req@ncols\@ne\fi
    \def\req@a{P}\ifx\req@t\req@a\advance\req@ncols\@ne\fi
    \def\req@a{M}\ifx\req@t\req@a\advance\req@ncols\@ne\fi}}

% ■見出し
%   直後が \begin{...}（表）のときは見出しを表の先頭行に組み込む（見出しだけがページ末に取り残されない）。
%   表が続かないとき（段落が続くとき）は単独の見出しとして出す。
\def\req@label{}
\newcommand{\req@plainlabel}[1]{\par\Needspace{4\baselineskip}\medskip\noindent ■#1\par\nopagebreak\smallskip}
\newcommand{\reqlabel}[1]{\@ifnextchar\begin{\gdef\req@label{#1}}{\req@plainlabel{#1}}}
\newcommand{\req@labelrow}{\ifx\req@label\@empty\else
  \multicolumn{\the\req@ncols}{@{}l@{}}{■\req@label}\\[.6ex]\fi}

% 本文の「・」箇条書き
\newenvironment{reqbullet}{%
  \begin{list}{・}{\setlength{\leftmargin}{1.5zw}\setlength{\labelwidth}{1zw}\setlength{\labelsep}{.5zw}%
    \setlength{\itemsep}{0pt}\setlength{\parsep}{0pt}\setlength{\topsep}{.3\baselineskip}}%
}{\end{list}}

% 画像（ファイルが無い間はテンプレート単体でコンパイルが通るよう枠を表示する）
\newcommand{\reqimage}[2][.9\textwidth]{%
  \IfFileExists{#2}%
    {\includegraphics[width=#1]{#2}}%
    {\fbox{\parbox[c][5cm][c]{\dimexpr#1-2\fboxsep-2\fboxrule\relax}{\centering ※図を添付（\detokenize{#2}）}}}}

% 表 : 列指定。比率は表幅に対する割合（セル余白と罫線ぶんを差し引いて表幅ぴったりに収める）
%   K・P列は左揃え（\raggedright）。狭い列で両端揃えにすると字間が不自然に空くため
\newcommand{\reqth}[1]{\textcolor{white}{\bfseries #1}}
\newcolumntype{K}[1]{|>{\bfseries\gtfamily\raggedright\arraybackslash}p{\dimexpr#1\dimexpr\linewidth-\arrayrulewidth\relax-2\tabcolsep-\arrayrulewidth\relax}}
\newcolumntype{P}[1]{|>{\raggedright\arraybackslash}p{\dimexpr#1\dimexpr\linewidth-\arrayrulewidth\relax-2\tabcolsep-\arrayrulewidth\relax}}
\newcolumntype{M}[1]{|>{\centering\arraybackslash}p{\dimexpr#1\dimexpr\linewidth-\arrayrulewidth\relax-2\tabcolsep-\arrayrulewidth\relax}}

% ヘッダー付き表（最初のデータ行が薄水色の縞模様）。直前の \reqlabel は先頭行として組み込まれる
\newenvironment{reqtable}[2]{%
  \par\smallskip\small
  \renewcommand{\arraystretch}{1.25}%
  \req@countcols{#1}%
  \arrayrulecolor{reqrule}%
  \ifx\req@label\@empty\rowcolors{2}{reqlight}{white}\else\rowcolors{2}{white}{reqlight}\fi% 先頭用・継続用ヘッダーの2行ぶんを数えたうえで，最初のデータ行が水色になる向き
  \begin{longtable}{#1|}
  \req@labelrow
  \hline
  \rowcolor{reqhead}#2\\
  \hline
  \endfirsthead
  \hline
  \rowcolor{reqhead}#2\\
  \hline
  \endhead
}{%
  \end{longtable}\gdef\req@label{}\par\smallskip}
% ヘッダー無しの表（項目名＋値）
\newenvironment{reqlist}[1]{%
  \par\smallskip\small
  \renewcommand{\arraystretch}{1.25}%
  \req@countcols{#1}%
  \arrayrulecolor{reqrule}%
  \ifx\req@label\@empty\rowcolors{1}{reqlight}{white}\else\rowcolors{2}{white}{reqlight}\fi
  \begin{longtable}{#1|}
  \req@labelrow
  \hline
  \endfirsthead
  \hline
  \endhead
}{%
  \end{longtable}\gdef\req@label{}\par\smallskip}
% よく使う表のプリセット
\newenvironment{reqcategory}{%
  \begin{reqtable}{K{.3}P{.7}}{\reqth{カテゴリー} & \reqth{詳細}}}{\end{reqtable}}
\newenvironment{reqitems}{%
  \begin{reqtable}{K{.25}P{.38}P{.37}}{\reqth{項目} & \reqth{説明} & \reqth{具体的な要件}}}{\end{reqtable}}
\makeatother


\begin{document}

%==============================================================================
% 表紙
%==============================================================================
\begin{titlepage}
\begin{center}
\vspace*{6em}
{\LARGE －【文書の副題】－}% サンプル: －要件定義書サンプル－

\vspace{6em}
\begingroup\setlength{\fboxrule}{3pt}\setlength{\fboxsep}{0pt}%
\fcolorbox{reqblue}{white}{\parbox[c][.26\textheight][c]{.72\textwidth}{\centering
{\huge\headfont 【システム名】}\\[.6em]
{\huge\headfont 要件定義書}}}
\endgroup

\vspace{9em}
{\Large 【YYYY年mm月dd日】}

\vspace{1.5em}
{\Large 【組織名・作成者】}
\end{center}

% 学校提出用に所属・氏名などが必要な場合は，スキル標準の表紙の書式で以下を追加する
%\vspace{4em}
%\begin{LARGE}
%\begin{flushright}
%【学年】情報工学科 【番号】番 \linebreak \vspace{1em}
%【氏名】
%\end{flushright}
%\end{LARGE}
%\vspace{2em}
%\begin{Large}
%\begin{tabular}{ll}
%提出期限 : & 【YYYY年MM月DD日（曜）hh:mm】 \\
%提出日 : & 【YYYY年MM月DD日（曜）hh:mm】
%\end{tabular}
%\end{Large}
\end{titlepage}
\setcounter{page}{2}% サンプルPDFに合わせ，表紙を1ページ目として数える（目次のページ番号がサンプルと同じ並びになる）

%==============================================================================
% 改訂履歴（サンプルp.2）
%==============================================================================
\section*{改訂履歴}
\begin{reqtable}{P{.18}P{.22}P{.28}P{.32}}{\reqth{バージョン} & \reqth{日付} & \reqth{改訂者} & \reqth{改訂内容}}
【1.0】 & 【日付】 & 【組織名】\par 氏名：【氏名】 & 【新規作成】 \\ \hline
 & & & \\ \hline
 & & & \\ \hline
 & & & \\ \hline
 & & & \\ \hline
\end{reqtable}

%==============================================================================
% 目次（サンプルp.3-4）
%==============================================================================
\tableofcontents

%==============================================================================
% 1. 概要（サンプルp.5-6）
%==============================================================================
\section{概要}

\subsection{システム化の目的}
\begin{reqbullet}
  \item 【システム化の目的1】
  \item 【システム化の目的2】
  \item 【システム化の目的3】
\end{reqbullet}

\subsection{現状の課題}
\begin{reqcategory}
【課題の分類1】 & ・【課題の詳細1】\par・【課題の詳細2】 \\ \hline
【課題の分類2】 & ・【課題の詳細1】\par・【課題の詳細2】 \\ \hline
【課題の分類3】 & ・【課題の詳細1】\par・【課題の詳細2】 \\ \hline
\end{reqcategory}

\subsection{プロジェクトの背景}
\begin{reqcategory}
【背景の分類1】 & ・【背景の詳細1】\par・【背景の詳細2】 \\ \hline
【背景の分類2】 & ・【背景の詳細1】\par・【背景の詳細2】 \\ \hline
【背景の分類3】 & ・【背景の詳細1】\par・【背景の詳細2】 \\ \hline
\end{reqcategory}

\subsection{システム構成図}
【システム名】の全体構成図を図\ref{fig:system-architecture}に示す。

\begin{figure}[H]
  \begin{center}
    \reqimage{assets/system-architecture.png}
    \caption{【システム名】の全体構成図}
    \label{fig:system-architecture}
  \end{center}
\end{figure}

\subsection{用語定義}
\begin{reqcategory}
【用語1】 & 【用語1の定義】 \\ \hline
【用語2】 & 【用語2の定義】 \\ \hline
【用語3】 & 【用語3の定義】 \\ \hline
\end{reqcategory}

%==============================================================================
% 2. 業務要件（サンプルp.7-11）
%==============================================================================
\section{業務要件}

\subsection{業務要件一覧}
本【システム名】は、以下の業務を実現するものとする。

\reqlabel{顧客対応業務}% 業務区分ごとに ■ブロックを増減する（サンプル: 顧客対応業務／内部業務／システム管理業務）
\begin{reqtable}{K{.22}P{.34}P{.24}M{.20}}{\reqth{業務} & \reqth{詳細} & \reqth{実施者} & \reqth{システム化対象}}
【業務名1】 & 【業務の詳細】 & 【実施者1】\par【実施者2】 & ○ \\ \hline
【業務名2】 & 【業務の詳細】 & 【実施者】 & \\ \hline% システム化対象外は空欄
\end{reqtable}

\reqlabel{内部業務}
\begin{reqtable}{K{.22}P{.34}P{.24}M{.20}}{\reqth{業務} & \reqth{詳細} & \reqth{実施者} & \reqth{システム化対象}}
【業務名1】 & 【業務の詳細】 & 【実施者】 & ○ \\ \hline
\end{reqtable}

\reqlabel{システム管理業務}
\begin{reqtable}{K{.22}P{.34}P{.24}M{.20}}{\reqth{業務} & \reqth{詳細} & \reqth{実施者} & \reqth{システム化対象}}
【業務名1】 & 【業務の詳細】 & 【実施者】 & \\ \hline
\end{reqtable}

\subsection{業務フロー図}
% サンプルの注記: 赤枠で囲むなどして、システム化範囲を明確にすること
\begin{figure}[H]
  \begin{center}
    \reqimage{assets/business-flow.png}
    \caption{業務フロー図}
    \label{fig:business-flow}
  \end{center}
\end{figure}

\subsection{規模}
\reqlabel{データ件数・サービス利用者数}
\begin{reqlist}{K{.4}P{.6}}
【項目1】 & 【値・説明】 \\ \hline
【項目2】 & 【値・説明】 \\ \hline
【項目3】 & 【値・説明】 \\ \hline
\end{reqlist}

\reqlabel{【取り扱い対象の規模】}% サンプル: 商品・サービスの取り扱い
\begin{reqlist}{K{.4}P{.6}}
【項目1】 & 【値・説明】 \\ \hline
【項目2】 & 【値・説明】 \\ \hline
\end{reqlist}

\reqlabel{【業務固有の規模】}% サンプル: オペレーターのエスカレーション処理
\begin{reqlist}{K{.4}P{.6}}
【項目1】 & 【値・説明】 \\ \hline
【項目2】 & 【値・説明】 \\ \hline
\end{reqlist}

\reqlabel{システムの規模}
\begin{reqlist}{K{.4}P{.6}}
【項目1】 & 【値・説明】 \\ \hline
【項目2】 & 【値・説明】 \\ \hline
\end{reqlist}

\subsection{場所}
【システムを利用する場所（例：社内・社外）】

\subsection{時期・時間}
\begin{reqlist}{K{.4}P{.6}}
システム稼働時間 & 【稼働時間】 \\ \hline
メンテナンス & 【メンテナンスの時期・時間】 \\ \hline
アクセスのピーク & 【ピーク時間帯】 \\ \hline
【その他の項目】 & 【値・説明】 \\ \hline
\end{reqlist}

\subsection{指標}
\reqlabel{パフォーマンス指標}
\begin{reqlist}{K{.4}P{.6}}
【指標名】 & 【目標値・説明】 \\ \hline
\end{reqlist}

\reqlabel{利用状況指標}
\begin{reqlist}{K{.4}P{.6}}
【指標名】 & 【目標値・説明】 \\ \hline
\end{reqlist}

\reqlabel{品質管理指標}
\begin{reqlist}{K{.4}P{.6}}
【指標名】 & 【目標値・説明】 \\ \hline
\end{reqlist}

\reqlabel{セキュリティ指標}
\begin{reqlist}{K{.4}P{.6}}
【指標名】 & 【目標値・説明】 \\ \hline
\end{reqlist}

\reqlabel{拡張性指標}
\begin{reqlist}{K{.4}P{.6}}
【指標名】 & 【目標値・説明】 \\ \hline
\end{reqlist}

\subsection{システム化範囲}
\begin{figure}[H]
  \begin{center}
    \reqimage{assets/system-scope.png}
    \caption{システム化範囲}
    \label{fig:system-scope}
  \end{center}
\end{figure}

\reqlabel{システム化の対象範囲}
\begin{reqtable}{K{.3}P{.7}}{\reqth{業務} & \reqth{内容}}
【業務区分1】 & ○【業務名1】\par○【業務名2】 \\ \hline
【業務区分2】 & ○【業務名1】\par○【業務名2】 \\ \hline
\end{reqtable}

\reqlabel{システム化の対象外範囲}
\begin{reqtable}{K{.3}P{.7}}{\reqth{業務} & \reqth{内容}}
【業務区分1】 & ○【業務名1】\par○【業務名2】 \\ \hline
\end{reqtable}

%==============================================================================
% 3. 機能要件（サンプルp.12-17）
%==============================================================================
\section{機能要件}

\subsection{機能一覧}
\reqlabel{基本機能}
\begin{reqtable}{K{.25}P{.40}P{.35}}{\reqth{機能} & \reqth{説明} & \reqth{対応業務要件}}
【機能名1】 & 【機能の説明】 & ○【対応する業務要件1】\par○【対応する業務要件2】 \\ \hline
【機能名2】 & 【機能の説明】 & ○【対応する業務要件】 \\ \hline
\end{reqtable}

\reqlabel{拡張機能}
\begin{reqtable}{K{.25}P{.40}P{.35}}{\reqth{機能} & \reqth{説明} & \reqth{対応業務要件}}
【機能名1】 & 【機能の説明】 & ○【対応する業務要件】 \\ \hline
\end{reqtable}

\subsection{画面}
\reqlabel{【利用者区分1】向け画面}% サンプル: 顧客向け画面／オペレーター向け画面／管理者向け画面
\begin{reqtable}{K{.3}P{.7}}{\reqth{画面名} & \reqth{説明}}
【画面名1】 & 【画面の説明】 \\ \hline
【画面名2】 & 【画面の説明】 \\ \hline
\end{reqtable}

\reqlabel{【利用者区分2】向け画面}
\begin{reqtable}{K{.3}P{.7}}{\reqth{画面名} & \reqth{説明}}
【画面名1】 & 【画面の説明】 \\ \hline
\end{reqtable}

\reqlabel{【利用者区分3】向け画面}
\begin{reqtable}{K{.3}P{.7}}{\reqth{画面名} & \reqth{説明}}
【画面名1】 & 【画面の説明】 \\ \hline
\end{reqtable}

\subsection{帳票}
\reqlabel{【帳票区分1】}% サンプル: レポート帳票／顧客向け帳票／管理者向け帳票
\begin{reqtable}{K{.3}P{.7}}{\reqth{帳票名} & \reqth{内容}}
【帳票名1】 & 【帳票の内容】 \\ \hline
\end{reqtable}

\reqlabel{【帳票区分2】}
\begin{reqtable}{K{.3}P{.7}}{\reqth{帳票名} & \reqth{内容}}
【帳票名1】 & 【帳票の内容】 \\ \hline
\end{reqtable}

\reqlabel{【帳票区分3】}
\begin{reqtable}{K{.3}P{.7}}{\reqth{帳票名} & \reqth{内容}}
【帳票名1】 & 【帳票の内容】 \\ \hline
\end{reqtable}

\subsection{情報・データ}
% サンプルの注記: 必要に応じて ER 図を添付
\begin{figure}[H]
  \begin{center}
    \reqimage{assets/er-diagram.png}
    \caption{ER図}
    \label{fig:er-diagram}
  \end{center}
\end{figure}

\begin{reqtable}{K{.3}P{.7}}{\reqth{情報・データ名} & \reqth{内容}}
【データ名1】 & 【データの内容】 \\ \hline
【データ名2】 & 【データの内容】 \\ \hline
\end{reqtable}

\subsection{外部インターフェース}
\reqlabel{【連携先の分類1】}% サンプル: 顧客対応関連／ナレッジベース関連／データ分析関連
\begin{reqtable}{K{.20}P{.30}P{.20}P{.30}}{\reqth{項目} & \reqth{説明} & \reqth{必要性} & \reqth{データ交換の仕様}}
【連携先1】 & 【インターフェースの説明】 & 【必要な理由】 & 【方式・認証・データ形式】 \\ \hline
\end{reqtable}

\reqlabel{【連携先の分類2】}
\begin{reqtable}{K{.20}P{.30}P{.20}P{.30}}{\reqth{項目} & \reqth{説明} & \reqth{必要性} & \reqth{データ交換の仕様}}
【連携先1】 & 【インターフェースの説明】 & 【必要な理由】 & 【方式・認証・データ形式】 \\ \hline
\end{reqtable}

\reqlabel{セキュリティ関連}
\begin{reqtable}{K{.20}P{.30}P{.20}P{.30}}{\reqth{項目} & \reqth{説明} & \reqth{必要性} & \reqth{プロトコル・認証方法}}
【項目1】 & 【インターフェースの説明】 & 【必要な理由】 & 【プロトコル・認証方法】 \\ \hline
\end{reqtable}

%==============================================================================
% 4. 非機能要件（サンプルp.18-39）
%   各表は「項目 | 説明 | 具体的な要件」(reqitems)．■ブロックは必要なものだけ残す
%==============================================================================
\section{非機能要件}

\subsection{ユーザビリティおよびアクセシビリティ}
\reqlabel{ユーザビリティ要件}
\begin{reqitems}
【項目】 & 【説明】 & ・【具体的な要件1】\par・【具体的な要件2】 \\ \hline
\end{reqitems}

\reqlabel{アクセシビリティ要件}
\begin{reqitems}
【項目】 & 【説明】 & ・【具体的な要件1】\par・【具体的な要件2】 \\ \hline
\end{reqitems}

\subsection{システム方式}
\reqlabel{システムアーキテクチャ}
\begin{reqitems}
【項目】 & 【説明】 & ・【具体的な要件1】\par・【具体的な要件2】 \\ \hline
\end{reqitems}

\reqlabel{データベース}
\begin{reqitems}
【項目】 & 【説明】 & ・【具体的な要件1】\par・【具体的な要件2】 \\ \hline
\end{reqitems}

\reqlabel{ネットワーク}
\begin{reqitems}
【項目】 & 【説明】 & ・【具体的な要件1】\par・【具体的な要件2】 \\ \hline
\end{reqitems}

\reqlabel{セキュリティ}
\begin{reqitems}
【項目】 & 【説明】 & ・【具体的な要件1】\par・【具体的な要件2】 \\ \hline
\end{reqitems}

\reqlabel{バックアップとリカバリ}
\begin{reqitems}
【項目】 & 【説明】 & ・【具体的な要件1】\par・【具体的な要件2】 \\ \hline
\end{reqitems}

\reqlabel{モニタリングとロギング}
\begin{reqitems}
【項目】 & 【説明】 & ・【具体的な要件1】\par・【具体的な要件2】 \\ \hline
\end{reqitems}

\subsection{規模}
\reqlabel{ユーザー数}
\begin{reqitems}
【項目】 & 【説明】 & 【具体的な要件】 \\ \hline
\end{reqitems}

\reqlabel{データ件数}
\begin{reqitems}
【項目】 & 【説明】 & 【具体的な要件】 \\ \hline
\end{reqitems}

\reqlabel{トランザクション数}
\begin{reqitems}
【項目】 & 【説明】 & 【具体的な要件】 \\ \hline
\end{reqitems}

\subsection{性能}
\reqlabel{レスポンスタイム（応答時間）}
\begin{reqitems}
【項目】 & 【説明】 & ・【具体的な要件1】\par・【具体的な要件2】 \\ \hline
\end{reqitems}

\reqlabel{スループット}
\begin{reqitems}
【項目】 & 【説明】 & 【具体的な要件】 \\ \hline
\end{reqitems}

\reqlabel{バッチ処理}
\begin{reqitems}
【項目】 & 【説明】 & ・【具体的な要件1】\par・【具体的な要件2】 \\ \hline
\end{reqitems}

\reqlabel{レイテンシー}
\begin{reqitems}
【項目】 & 【説明】 & ・【具体的な要件1】\par・【具体的な要件2】 \\ \hline
\end{reqitems}

\reqlabel{スケーラビリティ}
\begin{reqitems}
【項目】 & 【説明】 & ・【具体的な要件1】\par・【具体的な要件2】 \\ \hline
\end{reqitems}

\subsection{信頼性}
\reqlabel{可用性}
\begin{reqitems}
【項目】 & 【説明】 & ・【具体的な要件1】\par・【具体的な要件2】 \\ \hline
\end{reqitems}

\reqlabel{障害対応}
\begin{reqitems}
【項目】 & 【説明】 & 【具体的な要件】 \\ \hline
\end{reqitems}

\reqlabel{フェイルオーバー}
\begin{reqitems}
【項目】 & 【説明】 & ・【具体的な要件1】\par・【具体的な要件2】 \\ \hline
\end{reqitems}

\reqlabel{データの一貫性と耐久性}
\begin{reqitems}
【項目】 & 【説明】 & ・【具体的な要件1】\par・【具体的な要件2】 \\ \hline
\end{reqitems}

\reqlabel{モニタリングとアラート}
\begin{reqitems}
【項目】 & 【説明】 & ・【具体的な要件1】\par・【具体的な要件2】 \\ \hline
\end{reqitems}

\reqlabel{定期メンテナンス}
\begin{reqitems}
【項目】 & 【説明】 & ・【具体的な要件1】\par・【具体的な要件2】 \\ \hline
\end{reqitems}

\reqlabel{冗長化}
\begin{reqitems}
【項目】 & 【説明】 & ・【具体的な要件1】\par・【具体的な要件2】 \\ \hline
\end{reqitems}

\reqlabel{バックアップとリカバリ}
\begin{reqitems}
【項目】 & 【説明】 & ・【具体的な要件1】\par・【具体的な要件2】 \\ \hline
\end{reqitems}

\subsection{拡張性}
\begin{reqitems}
【項目】 & 【説明】 & ・【具体的な要件1】\par・【具体的な要件2】 \\ \hline
\end{reqitems}

\subsection{互換性}
\begin{reqitems}
【項目】 & 【説明】 & ・【具体的な要件1】\par・【具体的な要件2】 \\ \hline
\end{reqitems}

\subsection{継続性}
\reqlabel{冗長化}
\begin{reqitems}
【項目】 & 【説明】 & ・【具体的な要件1】\par・【具体的な要件2】 \\ \hline
\end{reqitems}

\reqlabel{バックアップとリカバリ}
\begin{reqitems}
【項目】 & 【説明】 & ・【具体的な要件1】\par・【具体的な要件2】 \\ \hline
\end{reqitems}

\reqlabel{サービスレベルアグリーメント（SLA）}
\begin{reqitems}
【項目】 & 【説明】 & ・【具体的な要件1】\par・【具体的な要件2】 \\ \hline
\end{reqitems}

\reqlabel{ビジネス継続計画（BCP）}
\begin{reqitems}
【項目】 & 【説明】 & ・【具体的な要件1】\par・【具体的な要件2】 \\ \hline
\end{reqitems}

\subsection{セキュリティ}
\reqlabel{データ保護}
\begin{reqitems}
【項目】 & 【説明】 & ・【具体的な要件1】\par・【具体的な要件2】 \\ \hline
\end{reqitems}

\reqlabel{認証と認可}
\begin{reqitems}
【項目】 & 【説明】 & ・【具体的な要件1】\par・【具体的な要件2】 \\ \hline
\end{reqitems}

\reqlabel{ネットワークセキュリティ}
\begin{reqitems}
【項目】 & 【説明】 & ・【具体的な要件1】\par・【具体的な要件2】 \\ \hline
\end{reqitems}

\reqlabel{脆弱性管理}
\begin{reqitems}
【項目】 & 【説明】 & ・【具体的な要件1】\par・【具体的な要件2】 \\ \hline
\end{reqitems}

\subsection{稼働環境}
\begin{reqitems}
【項目】 & 【説明】 & ・【具体的な要件1】\par・【具体的な要件2】 \\ \hline
\end{reqitems}

\subsection{テスト}
\begin{reqitems}
【項目】 & 【説明】 & ・【具体的な要件1】\par・【具体的な要件2】 \\ \hline
\end{reqitems}

\subsection{移行}
% サンプルの注記: システムリプレイス時など、必要な場合のみ記載する（不要ならこの節ごと削除）
\reqlabel{データ移行}
\begin{reqitems}
【項目】 & 【説明】 & ・【具体的な要件1】\par・【具体的な要件2】 \\ \hline
\end{reqitems}

\reqlabel{移行プロセス}
\begin{reqitems}
【項目】 & 【説明】 & ・【具体的な要件1】\par・【具体的な要件2】 \\ \hline
\end{reqitems}

\reqlabel{移行テスト}
\begin{reqitems}
【項目】 & 【説明】 & ・【具体的な要件1】\par・【具体的な要件2】 \\ \hline
\end{reqitems}

\reqlabel{移行後サポート}
\begin{reqitems}
【項目】 & 【説明】 & ・【具体的な要件1】\par・【具体的な要件2】 \\ \hline
\end{reqitems}

\reqlabel{リスク管理}
\begin{reqitems}
【項目】 & 【説明】 & ・【具体的な要件1】\par・【具体的な要件2】 \\ \hline
\end{reqitems}

\subsection{引継ぎ}
\reqlabel{ドキュメント作成}
\begin{reqitems}
【項目】 & 【説明】 & ・【具体的な要件1】\par・【具体的な要件2】 \\ \hline
\end{reqitems}

\reqlabel{トレーニングと教育}
\begin{reqitems}
【項目】 & 【説明】 & ・【具体的な要件1】\par・【具体的な要件2】 \\ \hline
\end{reqitems}

\reqlabel{サポート体制}
\begin{reqitems}
【項目】 & 【説明】 & ・【具体的な要件1】\par・【具体的な要件2】 \\ \hline
\end{reqitems}

\reqlabel{コミュニケーション}
\begin{reqitems}
【項目】 & 【説明】 & ・【具体的な要件1】\par・【具体的な要件2】 \\ \hline
\end{reqitems}

\subsection{教育}
\reqlabel{ユーザートレーニング}
\begin{reqitems}
【項目】 & 【説明】 & ・【具体的な要件1】\par・【具体的な要件2】 \\ \hline
\end{reqitems}

\reqlabel{トレーニング教材}
\begin{reqitems}
【項目】 & 【説明】 & ・【具体的な要件1】\par・【具体的な要件2】 \\ \hline
\end{reqitems}

\reqlabel{トレーニングプログラムの実施}
\begin{reqitems}
【項目】 & 【説明】 & ・【具体的な要件1】\par・【具体的な要件2】 \\ \hline
\end{reqitems}

\subsection{運用}
\reqlabel{運用体制}
\begin{reqitems}
【項目】 & 【説明】 & ・【具体的な要件1】\par・【具体的な要件2】 \\ \hline
\end{reqitems}

\reqlabel{運用手順}
\begin{reqitems}
【項目】 & 【説明】 & ・【具体的な要件1】\par・【具体的な要件2】 \\ \hline
\end{reqitems}

\reqlabel{モニタリングとアラート}
\begin{reqitems}
【項目】 & 【説明】 & ・【具体的な要件1】\par・【具体的な要件2】 \\ \hline
\end{reqitems}

\reqlabel{バックアップとリカバリ}
\begin{reqitems}
【項目】 & 【説明】 & ・【具体的な要件1】\par・【具体的な要件2】 \\ \hline
\end{reqitems}

\reqlabel{定期メンテナンス}
\begin{reqitems}
【項目】 & 【説明】 & ・【具体的な要件1】\par・【具体的な要件2】 \\ \hline
\end{reqitems}

\reqlabel{運用レポート}
\begin{reqitems}
【項目】 & 【説明】 & ・【具体的な要件1】\par・【具体的な要件2】 \\ \hline
\end{reqitems}

\reqlabel{運用ツール}
\begin{reqitems}
【項目】 & 【説明】 & ・【具体的な要件1】\par・【具体的な要件2】 \\ \hline
\end{reqitems}

\reqlabel{コミュニケーション}
\begin{reqitems}
【項目】 & 【説明】 & ・【具体的な要件1】\par・【具体的な要件2】 \\ \hline
\end{reqitems}

\subsection{保守}
\reqlabel{保守体制}
\begin{reqitems}
【項目】 & 【説明】 & ・【具体的な要件1】\par・【具体的な要件2】 \\ \hline
\end{reqitems}

\reqlabel{保守手順}
\begin{reqitems}
【項目】 & 【説明】 & ・【具体的な要件1】\par・【具体的な要件2】 \\ \hline
\end{reqitems}

\reqlabel{障害対応}
\begin{reqitems}
【項目】 & 【説明】 & ・【具体的な要件1】\par・【具体的な要件2】 \\ \hline
\end{reqitems}

\reqlabel{ハードウェア保守}
\begin{reqitems}
【項目】 & 【説明】 & ・【具体的な要件1】\par・【具体的な要件2】 \\ \hline
\end{reqitems}

\reqlabel{ソフトウェア保守}
\begin{reqitems}
【項目】 & 【説明】 & ・【具体的な要件1】\par・【具体的な要件2】 \\ \hline
\end{reqitems}

\reqlabel{保守記録}
\begin{reqitems}
【項目】 & 【説明】 & ・【具体的な要件1】\par・【具体的な要件2】 \\ \hline
\end{reqitems}

\reqlabel{外部ベンダー管理}
\begin{reqitems}
【項目】 & 【説明】 & ・【具体的な要件1】\par・【具体的な要件2】 \\ \hline
\end{reqitems}

\reqlabel{保守ツール}
\begin{reqitems}
【項目】 & 【説明】 & ・【具体的な要件1】\par・【具体的な要件2】 \\ \hline
\end{reqitems}

%==============================================================================
% 5. 補足（サンプルp.40）
%==============================================================================
\section{補足}

\subsection{プロジェクト　スケジュール}
\begin{figure}[H]
  \begin{center}
    \reqimage{assets/project-schedule.png}
    \caption{プロジェクトスケジュール}
    \label{fig:project-schedule}
  \end{center}
\end{figure}

\subsection{プロジェクト　体制}
\begin{figure}[H]
  \begin{center}
    \reqimage{assets/project-organization.png}
    \caption{プロジェクト体制}
    \label{fig:project-organization}
  \end{center}
\end{figure}

\begin{flushright}
以上
\end{flushright}

% 参考文献が必要な場合は，.bib を用意して以下を有効にする（\cite{キー} で本文から引用）
%\bibliographystyle{jplain}
%\bibliography{requirements-template} % .bib拡張子を除いたファイル名

\end{document}
